\SourceLeaf{186}%
% A544-P0186-D01
das gezählte ein Entgegengesetztes hat, was mit ihm vereinigt gedacht der Vernichtung gleich zu stellen ist. Genau besehen findet diese Voraussetzung nur da Statt, wo nicht Substanzen (für sich denkbare Gegenstände) sondern Relationen zwischen je zweien Gegenständen das gezählte sind. Postulirt wird dabei, dass diese Gegenstände auf eine bestimmte Art in eine Reihe geordnet sind z. B. $A,B,C,D\dotsfour$, und dass die Relation des $A$ zu $B$ als der Relation des $B$ zu $C$ u. s. w. gleich betrachtet werden kann. Hier gehört nun zu dem Begriff der Entgegensetzung nichts weiter als der \emph{Umtausch} der Glieder der Relation, so dass wenn die Relation (oder der Uebergang) von $A$ zu $B$ als $+1$ gilt, die Relation von $B$ zu $A$ durch $-1$ dargestellt werden muss. Insofern also eine solche Reihe auf beiden Seiten unbegrenzt ist, repräsentirt jede reelle ganze Zahl die Relation eines beliebig als Anfang gewählten Gliedes zu einem bestimmten Gliede der Reihe.

% A544-P0186-D02
Sind aber die Gegenstände von solcher Art, dass sie nicht in Eine, wenn gleich unbegrenzte, Reihe geordnet werden können, sondern sich nur in Reihen von Reihen ordnen lassen, oder was dasselbe ist, bilden sie eine Mannigfaltigkeit von zwei Dimensionen; verhält es sich dann mit den Relationen einer Reihe zu einer andern oder den Uebergängen aus einer in die andere auf eine ähnliche Weise wie vorhin mit den Uebergängen von einem Gliede einer Reihe zu einem andern Gliede derselben Reihe, so bedarf es offenbar zur Abmessung des Ueberganges von einem Gliede des Systems zu einem andern ausser den vorigen Einheiten $+1$ und $-1$ noch zweier andern unter sich auch entgegengesetzten $+i$ und $-i$. Offenbar muss aber dabei noch postulirt werden, dass die Einheit $i$ allemal den Uebergang von einem gegebenen Gliede einer Reihe zu einem \emph{bestimmten} Gliede der unmittelbar angrenzenden Reihe bezeichne. Auf diese Weise wird also das System auf eine doppelte Art in Reihen von Reihen geordnet werden können.

% A544-P0186-D03
Der Mathematiker abstrahirt gänzlich von der Beschaffenheit der Gegenstände und dem Inhalt ihrer Relationen; er hat es blos mit der Abzählung und Vergleichung der Relationen unter sich zu thun: insofern ist er eben so, wie er den durch $+1$ und $-1$ bezeichneten Relationen, an sich betrachtet, Gleichartigkeit beilegt, solche auf alle vier Elemente $+1$, $-1$, $+i$ und $-i$ zu erstrecken befugt.

% A544-P0186-D04
Zur Anschauung lassen sich diese Verhältnisse nur durch eine Darstellung
\SourceLeaf{187}%
% A544-P0187-D01
im Raume bringen, und der einfachste Fall ist, wo kein Grund vorhanden ist, die Symbole der Gegenstände anders als quadratisch anzuordnen, indem man nemlich eine unbegrenzte Ebene durch zwei Systeme von Parallellinien, die einander rechtwinklig durchkreuzen, in Quadrate vertheilt, und die Durchschnittspunkte zu den Symbolen wählt. Jeder solche Punkt $A$ hat hier vier Nachbaren, und wenn man die Relation des $A$ zu einem benachbarten Punkte durch $+1$ bezeichnet, so ist die durch $-1$ zu bezeichnende von selbst bestimmt, während man, welche der beiden andern man will, für $+i$ wählen, oder den sich auf $+i$ beziehenden Punkt nach Gefallen \emph{rechts} oder \emph{links} nehmen kann. Dieser Unterschied zwischen rechts und links ist, so bald man vorwärts und rückwärts \emph{in} der Ebene, und oben und unten in Beziehung auf die beiden Seiten der Ebene einmal (nach Gefallen) festgesetzt hat, \emph{in sich} völlig bestimmt, wenn wir gleich unsere Anschauung dieses Unterschiedes andern \emph{nur} durch Nachweisung an wirklich vorhandenen materiellen Dingen mittheilen können\begingroup\renewcommand{\thefootnote}{*)}\footnote{Beide Bemerkungen hat schon \textsc{Kant} gemacht, aber man begreift nicht, wie dieser scharfsinnige Philosoph in der ersteren einen Beweis für seine Meinung, dass der Raum \emph{nur} Form unserer äussern Anschauung sei, zu finden glauben konnte, da die zweite so klar das Gegentheil, und dass der Raum unabhängig von unserer Anschauungsart eine reelle Bedeutung haben muss, beweiset.}\endgroup. Wenn man aber auch über letzteres sich entschlossen hat, sieht man, dass es doch von unserer Willkür abhing, welche von den beiden in Einem Punkte sich durchkreuzenden Reihen wir als Hauptreihe, und welche Richtung in ihr man als auf positive Zahlen sich beziehend ansehen wollten; man sieht ferner, dass wenn man die vorher als $+i$ behandelte Relation für $+1$ nehmen will, man nothwendig die vorher durch $-1$ bezeichnete Relation für $+i$ nehmen muss. Das heisst aber, in der Sprache der Mathematiker, $+i$ ist mittlere Proportionalgrösse zwischen $+1$ und $-1$ oder entspricht dem Zeichen $\surd-1$: wir sagen absichtlich nicht \emph{die} mittlere Proportionalgrösse, denn $-i$ hat offenbar gleichen Anspruch. Hier ist also die Nachweisbarkeit einer anschaulichen Bedeutung von $\surd-1$ vollkommen gerechtfertigt, und mehr bedarf es nicht, um diese Grösse in das Gebiet der Gegenstände der Arithmetik zuzulassen.

% A544-P0187-D02
Wir haben geglaubt, den Freunden der Mathematik durch diese kurze Darstellung der Hauptmomente einer neuen Theorie der sogenannten imaginären Grössen einen Dienst zu erweisen. Hat man diesen Gegenstand bisher aus einem falschen Gesichtspunkt betrachtet und eine geheimnissvolle Dunkelheit dabei 
\SourceLeaf{188}%
% A544-P0188-D01
gefunden, so ist dies grossentheils den wenig schicklichen Benennungen zuzuschreiben. Hätte man $+1$, $-1$, $\surd-1$ nicht positive, negative, imaginäre (oder gar unmögliche) Einheit, sondern etwa directe, inverse, laterale Einheit genannt, so hätte von einer solchen Dunkelheit kaum die Rede sein können. Der Verf. hat sich vorbehalten, den Gegenstand, welcher in der vorliegenden Abhandlung eigentlich nur gelegentlich berührt ist, künftig vollständiger zu bearbeiten, wo dann auch die Frage, warum die Relationen zwischen Dingen, die eine Mannigfaltigkeit von mehr als zwei Dimensionen darbieten, nicht noch andere in der allgemeinen Arithmetik zulässige Arten von Grössen liefern können, ihre Beantwortung finden wird.
\BookRule{.28\linewidth}

\clearpage
\SourceLeaf{189}%
% A544-P0189-D01
\thispagestyle{empty}
\pdfbookmark[0]{Anzeigen nicht eigner Schriften}{other-notices-title}
\begin{center}
\vspace*{1in}{\Large ANZEIGEN\par}
\vspace{.4in}{\large NICHT EIGNER\par}
\vspace{.4in}{\Large SCHRIFTEN.\par}
\end{center}

